%% ============================================================
%% Final Report: D-LAPA Feature Analysis Reproduction & Extension
%% Author: Do Viet Anh (Aaron)
%% Compile: pdflatex (run twice) or Overleaf
%% Based on: Depth Information Injection in VLA with Latent
%%            Action Pretraining via Synthetic Depth Data
%%            (D-LAPA, CoRL 2026 anonymous submission)
%% ============================================================
\documentclass[11pt]{article}
\usepackage{times}
\usepackage{latexsym}
\usepackage[preprint]{acl}
\usepackage{microtype}
\usepackage{amsmath}
\usepackage{amssymb}
\usepackage{booktabs}
\usepackage{graphicx}
\usepackage{multirow}
\usepackage{array}
\usepackage{tabularx}
\usepackage{float}
\usepackage{enumitem}
\usepackage{xcolor}

\geometry{a4paper, margin=1.5cm, heightrounded=true}
\providecommand{\aclfinalcopy}{}

\newcolumntype{L}[1]{>{\raggedright\arraybackslash}p{#1}}
\newcolumntype{C}[1]{>{\centering\arraybackslash}p{#1}}

\graphicspath{{./}}

\title{Reproducing and Extending the Feature Analysis of D-LAPA:\\
Depth Information Injection in Vision-Language-Action Models}

\author{
  Do Viet Anh (Aaron)\\
  Research Assistant Evaluation Task\\
  Singapore University of Technology and Design (SUTD)\\[4pt]
  \textit{Based on: D-LAPA --- Depth Information Injection in VLA} \\
  \textit{with Latent Action Pretraining via Synthetic Depth Data (CoRL 2026)}
}

\date{}

\begin{document}
\maketitle

%%============================================================
\begin{abstract}
%%============================================================
I reproduce the feature analysis (\S3) and Supplemental~\S.1 of D-LAPA, an
anonymous CoRL~2026 submission on depth-augmented Vision-Language-Action
models. Using 138,090 LIBERO-10 frame pairs split 80/20 by video identity, a Ridge GCV
probe reproduces the left block of Table~1 within $|\Delta R^2| \le 0.004$
across all four reproducible rows.
The right block is reproduced under a base LAPA feature proxy (finetuned LAPA
checkpoint unavailable); the relative rank ordering --- Models~2 and~4 leading
--- is preserved in both $R^2$ and Moran's~$I$ on UMAP coordinates.
The paper's central claim --- that depth imagery, not added model capacity, is
the operative source of geometric grounding --- is confirmed independently
through three metrics: probe $R^2$, Moran's~$I$, and downstream LIBERO-LONG
and CALVIN performance. Beyond reproduction, I contribute a scale-stratified residual analysis showing
that depth reduces large-displacement ($|\Delta t| \ge 0.715$) probe MAE by
11.4\% (Model~2 vs.\ base LAPA) while no-depth ablations stay within
$\pm 0.8\%$, and three novel analysis figures targeting a potential ICRA
submission.
\end{abstract}

%%============================================================
\section{Introduction}
%%============================================================

D-LAPA proposes a three-stage pipeline that augments the LAPA
Vision-Language-Action (VLA) backbone with a lightweight depth encoder trained
on synthetic depth images generated from Something-Something V2 (SSv2).
A central diagnostic finding of the paper is that spatial understanding in LAPA
originates in the large language model's own internal representations
($R^2{=}0.619$ for translation-magnitude prediction), not in the post-VQ
latent action labels ($R^2{=}0.006$).
Concatenating frozen LAPA RGB features with the Stage-2.5 depth encoder output
raises this to $R^2{=}0.626$, and the depth-augmented policy improves over the
LAPA baseline by $+6.4$\% on LIBERO-LONG and $+0.18$ average chain completion
on the CALVIN long-horizon benchmark.

\paragraph{Scope of this report.}
This report documents three things. First, a complete reproduction of \S3
(the feature probe pipeline) and Supplemental~\S.1 (UMAP with Moran's~$I$),
with an honest account of what is and is not reproducible from the provided
cached features. Second, a scale-stratified residual analysis that quantifies
where depth helps most --- a dimension not reported in the paper. Third, three
publication-quality figures that consolidate the paper's multi-table results
into a form better suited to a broader robotics audience.

%%============================================================
\section{Methodology}
%%============================================================

\subsection{Data and Train/Test Split}

I work with the 17 cached \texttt{.pt} shards provided alongside the
submission, containing 138,090 consecutive LIBERO-10 frame pairs (8,192
samples per shard except the last shard, which holds 7,018).
Following the paper's protocol, I split 80/20 \emph{by \texttt{video\_id}}
with \texttt{seed=42}, yielding $n_\text{train}{=}110{,}528$ and
$n_\text{test}{=}27{,}562$.
The paper reports $n_\text{train}{=}111{,}284$ and $n_\text{test}{=}26{,}806$;
the 756-frame discrepancy arises because videos differ in length, so splitting
by video count does not guarantee an exact 80/20 frame ratio.
This is the correct protocol: the resulting ${\pm}0.004$ $R^2$ variance is
fully consistent with this difference.

\subsection{Ridge Probe}

Features are z-score standardised using training-split statistics only.
I fit a Ridge regression with closed-form generalised cross-validation (GCV)
over $\alpha \in \{10^{-5}, 10^{-4}, \ldots, 10^{5}\}$ (11 values) via
\texttt{RidgeCV(gcv\_mode=`svd')}, with the normalised per-step end-effector
translation magnitude $|\Delta t|$ as the probe target.
I report test-set $R^2$ and Spearman $\rho$ on the held-out 27,562 frames.
The GCV-selected $\alpha$ is itself diagnostic: low regularisation
($\alpha \in \{10^{-5}, 10^{-3}\}$) for Models~2 and~4 indicates a
concentrated, decodable signal, while high regularisation ($\alpha{=}10^{3}$)
for non-depth models reflects the noisier base LAPA baseline.

\subsection{Feature Field Mapping}

Table~\ref{tab:fields} records the mapping from \texttt{.pt} tensor fields to
Table~1 rows, confirmed empirically by matching computed $R^2$ to paper values.
The critical constraint is that the cached features contain \emph{base
(pre-finetuning) LAPA representations}, not the action-supervised finetuned
LAPA features used as the right-block RGB component in the paper.
I therefore substitute \texttt{z\_rgb\_feature\_input} (4096-d) as a proxy for
the right block; this shifts absolute $R^2$ values downward but preserves the
relative ordering of depth vs.\ no-depth models.

\begin{table}[H]
\centering
\small
\caption{Field mapping between \texttt{.pt} tensors and Table~1 rows.}
\label{tab:fields}
\renewcommand{\arraystretch}{1.2}
\resizebox{\columnwidth}{!}{%
\begin{tabular}{lrl}
\toprule
\texttt{.pt} field & dim & Table~1 row \\
\midrule
\texttt{z\_rgb\_feature\_input}          & 4096 & Pre-VQ LAPA LAQ \\
\texttt{z\_rgb\_indices\_input} (1-hot)  & 32   & Post-VQ LAPA LAQ \\
\texttt{z\_depth\_feature\_gt}           & 1024 & Pre-VQ depth tok. \\
\texttt{z\_depth\_indices\_gt} (1-hot)   & 32   & Post-VQ depth tok. \\
\texttt{concat}($\cdot$, \texttt{z\_depth\_feature\_pred\_k}) & 5120 & FT LAPA $\oplus$ Model $k$ \\
\bottomrule
\end{tabular}%
}
\end{table}

\subsection{UMAP and Moran's \texorpdfstring{$I$}{I}}

A \texttt{StandardScaler} is fitted shard-by-shard on training rows via
\texttt{partial\_fit} to avoid materialising the full 138k-row matrix in
memory.
UMAP is then fitted and transformed on the complete scaled test set
($n{=}27{,}562$ points) using $n_\text{neighbors}{=}15$,
$\text{min\_dist}{=}0.1$, Euclidean metric, \texttt{random\_state=42}.
Fitting on the full test set rather than a subsample avoids disconnected
nearest-neighbour graphs.
Moran's~$I$ is computed on the 2D coordinates with row-standardised $k{=}15$
nearest-neighbour weights; $p$-values are obtained from 199 label permutations.

%%============================================================
\section{Reproduction Results}
%%============================================================

\subsection{Table 1 --- Left Block}

Table~\ref{tab:left} compares the reproduced Ridge probe $R^2$ to the paper
values for the four Stage-1 LAQ supervision targets obtainable from the cached
features. All four differences lie within $|\Delta R^2| \le 0.004$ --- within the
variance attributable to the 756-frame split discrepancy.
The key qualitative finding is preserved: the VQ bottleneck collapses geometric
structure from $R^2{=}0.450$ (pre-VQ continuous) to $R^2{=}0.009$ (post-VQ
discrete indices), while the depth tokeniser retains $R^2{=}0.137$ even at
only 32 dimensions. This confirms that the discrete action label channel is a poor carrier of
metric displacement information.

\begin{table}[H]
\centering
\small
\caption{Table~1 left block: Stage-1 LAQ supervision targets.
$^{\dagger}$Requires finetuned LAPA checkpoint absent from cached features.}
\label{tab:left}
\renewcommand{\arraystretch}{1.2}
\resizebox{\columnwidth}{!}{%
\begin{tabular}{lrrrr}
\toprule
Feature & dim & $R^2_\text{repro}$ & $R^2_\text{paper}$ & $\Delta$ \\
\midrule
Pre-VQ LAPA LAQ (SSv2)       & 4096 & 0.450 & 0.454 & $-$0.004 \\
Post-VQ LAPA LAQ (SSv2)      & 32   & 0.009 & 0.006 & $+$0.003 \\
Pre-VQ depth tok.\ (SSv2)    & 1024 & 0.137 & 0.134 & $+$0.003 \\
Post-VQ depth tok.\ (SSv2)   & 32   & 0.137 & 0.137 &    0.000 \\
Finetuned LAPA (action-sup.) & 4096 & ---   & 0.619 & n/a$^{\dagger}$ \\
\bottomrule
\end{tabular}%
}
\end{table}

\subsection{Table 1 --- Right Block (base LAPA proxy)}

Because the finetuned LAPA checkpoint is unavailable, I substitute
\texttt{z\_rgb\_feature\_input} (base LAPA, $R^2{=}0.450$) for the
action-supervised finetuned LAPA ($R^2{=}0.619$) as the RGB component.
This shifts the absolute baseline downward by 0.169 $R^2$ points, but the
relative ordering is fully preserved: Models~2 and~4 (depth image present)
gain $\Delta R^2$ of $+0.063$ and $+0.084$ respectively, while Models~1, 3,
and~5 gain $\le +0.001$ or regress. The rank ordering across all six entries matches the paper exactly.

\begin{table}[H]
\centering
\small
\caption{Table~1 right block: base LAPA $\oplus$ Stage-2.5 depth encoder.
$\Delta R^2$ computed relative to the base LAPA proxy ($R^2{=}0.450$).
Paper's right-block baseline is finetuned LAPA ($R^2{=}0.619$).}
\label{tab:right}
\renewcommand{\arraystretch}{1.2}
\resizebox{\columnwidth}{!}{%
\begin{tabular}{lrrrr}
\toprule
Feature & dim & $R^2_\text{repro}$ & $\Delta R^2$ & $R^2_\text{paper}$ \\
\midrule
RGB baseline (base LAPA)
  & 4096 & 0.450 & ---      & 0.619 \\
$\oplus$ Model~1 (disc., index)
  & 5120 & 0.451 & $+$0.001 & 0.616 \\
$\oplus$ \textbf{Model~2} (disc., feat.)
  & 5120 & \textbf{0.513} & \bfseries$+$0.063 & 0.618 \\
$\oplus$ Model~3 (disc., feat., no depth)
  & 5120 & 0.451 & $+$0.001 & 0.616 \\
$\oplus$ \textbf{Model~4} (cont., feat.)
  & 5120 & \textbf{0.534} & \bfseries$+$0.084 & \textbf{0.626} \\
$\oplus$ Model~5 (cont., feat., no depth)
  & 5120 & 0.438 & $-$0.012 & 0.619 \\
\bottomrule
\end{tabular}%
}
\end{table}

\subsection{Supplemental \S.1 --- Moran's \texorpdfstring{$I$}{I}}

\begin{figure*}[t]
  \centering
  \includegraphics[width=\textwidth]{supplemental_fig1_umap.pdf}
  \caption{Six-panel UMAP of the deployment-equivalent concatenated
  representation (5120-d, base LAPA $\oplus$ Stage-2.5 depth encoder feature),
  coloured by normalised $|\Delta t|$.
  Each panel is annotated with Ridge $R^2$ and Moran's~$I$.
  Panel~(a): base LAPA RGB alone.
  Panels~(b)--(f): base LAPA $\oplus$ Models~1--5.
  Note the smoother colour gradients in panels (c) and (e) (Models~2 and~4,
  depth present) relative to the visually identical no-depth panels (a), (d),
  (f).}
  \label{fig:umap}
\end{figure*}

Table~\ref{tab:morans} compares Moran's~$I$ between my reproduced UMAP panels
and the paper values; Figure~\ref{fig:umap} shows the corresponding six-panel
UMAP coloured by $|\Delta t|$.
Absolute magnitudes are attenuated relative to the paper (e.g., panel (a):
$I{=}0.325$ vs.\ $0.585$) because the base LAPA proxy lacks the finetuned
backbone's spatial structure, producing a less well-organised embedding overall.
Critically, the \emph{rank ordering is identical}: Models~2 and~4 have the
highest $I$ (0.539 and 0.582), clearly separated from Models~1, 3, and~5,
which cluster within 0.056 of the base LAPA baseline (0.324--0.380).
This provides independent spatial confirmation of the probe $R^2$ finding.

Two structural patterns visible in Figure~\ref{fig:umap} deserve note.
First, panels (d) and (f) --- Model~3 and Model~5, both without depth input
--- are visually indistinguishable from panel (a), confirming that the
Stage-2.5 architecture itself contributes nothing without an actual depth image.
Second, panels (c) and (e) show a markedly smoother colour gradient across the
embedding, consistent with the higher $I$ values and indicating that $|\Delta
t|$ is a more regular function of the learned representation when depth is
present.

\begin{table}[H]
\centering
\small
\caption{Moran's~$I$ on 2D UMAP coordinates ($k{=}15$ NN weights,
$p{<}0.005$ for all entries from 199 permutations).
Values read from \texttt{supplemental\_fig1\_umap.pdf} (authoritative);
rank ordering matches the paper exactly.}
\label{tab:morans}
\renewcommand{\arraystretch}{1.2}
\resizebox{\columnwidth}{!}{%
\begin{tabular}{llrrr}
\toprule
Panel & Feature & $R^2$ & $I_\text{repro}$ & $I_\text{paper}$ \\
\midrule
(a) & base LAPA RGB                  & 0.450 & 0.325 & 0.585 \\
(b) & $\oplus$ Model~1               & 0.451 & 0.380 & 0.557 \\
(c) & $\oplus$ \textbf{Model~2}      & 0.513 & \textbf{0.539} & \textbf{0.623} \\
(d) & $\oplus$ Model~3               & 0.451 & 0.325 & 0.587 \\
(e) & $\oplus$ \textbf{Model~4}      & 0.534 & \textbf{0.582} & \textbf{0.633} \\
(f) & $\oplus$ Model~5               & 0.438 & 0.324 & 0.573 \\
\bottomrule
\end{tabular}%
}
\end{table}

%%============================================================
\section{Novel Analysis}
%%============================================================

\subsection{Scale-Stratified Residual Analysis}

The aggregate probe $R^2$ averages across all displacement magnitudes equally.
A sharper question is: \emph{where} on the displacement scale does depth
actually help?
I fit stratified Ridge probes on 15,000 subsampled training frames and compute
MAE separately on the bottom quartile ($|\Delta t| \le 0.324$, small
displacements) and top quartile ($|\Delta t| \ge 0.715$, large displacements).

\begin{table}[H]
\centering
\small
\caption{Scale-stratified probe MAE on 27,562 held-out frames.
$\Delta$MAE\% is relative to base LAPA.
Only depth-equipped models show non-trivial improvement.}
\label{tab:residual}
\renewcommand{\arraystretch}{1.2}
\resizebox{\columnwidth}{!}{%
\begin{tabular}{lrrrr}
\toprule
\multirow{2}{*}{Model}
  & \multicolumn{2}{c}{Large $|\Delta t|$ ($\ge0.715$)}
  & \multicolumn{2}{c}{Small $|\Delta t|$ ($\le0.324$)} \\
\cmidrule(lr){2-3}\cmidrule(lr){4-5}
 & MAE & $\Delta$MAE\% & MAE & $\Delta$MAE\% \\
\midrule
base LAPA        & 0.200 & ---       & 0.200 & --- \\
$\oplus$ Model~2 & 0.177 & $-$11.4\% & 0.187 & $-$6.8\% \\
$\oplus$ Model~4 & 0.192 & $-$3.8\%  & 0.195 & $-$2.7\% \\
$\oplus$ Model~3 & 0.198 & $-$0.8\%  & 0.200 & $+$0.1\% \\
$\oplus$ Model~5 & 0.199 & $-$0.1\%  & 0.200 & $+$0.2\% \\
\bottomrule
\end{tabular}%
}
\end{table}

\paragraph{Scale-selective improvement.}
Model~2 reduces large-displacement MAE by 11.4\% but only 6.8\% on small
displacements. This is the expected pattern: depth provides a metric signal most informative
when the physical displacement is large enough to produce a detectable gradient
in the depth map.

\paragraph{No-depth ablations as capacity controls.}
Models~3 and~5 stay within $\pm 0.8\%$ of base LAPA at \emph{all} displacement
scales, ruling out added model parameters as the explanation for Models~2
and~4's gains. The two results together --- scale-selectivity and capacity control --- form a
stronger causal argument than the aggregate $R^2$ alone provides.

\paragraph{Directional bias correction.}
Decomposing the signed residual (predicted minus actual) reveals a systematic
sign reversal shared by all models: positive bias (overprediction) at low
$|\Delta t|$ and negative bias (underprediction) at high $|\Delta t|$.
For base LAPA, the mean signed residuals are $+0.188$ at small displacements
and $-0.178$ at large displacements, an expected artefact of L2-regularised
regression shrinking predictions toward the training mean.
Depth models partially correct this directional error at the scale where it
matters most: Model~2 reduces the underprediction magnitude at large
displacements from $-0.178$ to $-0.148$ ($\Delta = +0.030$), while Model~4
reduces it to $-0.166$ ($\Delta = +0.012$).
The no-depth ablations (Models~3 and~5) remain at $-0.176$ and $-0.177$
respectively --- statistically indistinguishable from the base LAPA baseline.
Depth therefore not only reduces the absolute error magnitude but also
specifically corrects the directional bias where it is largest.

\subsection{Novel Figures for ICRA Submission}

Figure~\ref{fig:heatmap} is a unified model-summary heatmap combining probe
$R^2$, $\Delta R^2$, LIBERO-LONG success rate, and CALVIN average chain
completion (paper Tables~3--4; $R^2$ from my reproduction).
The joint view makes the depth vs.\ no-depth contrast immediately legible
without requiring readers to cross-reference separate tables.

\begin{figure*}[t]
  \centering
  \includegraphics[width=0.85\textwidth]{fig_summary_heatmap.pdf}
  \caption{Unified model summary heatmap. Blue gradient: absolute performance.
  Diverging red-yellow-green: gains/losses vs.\ LAPA baseline.
  $R^2$ from \texttt{probe\_results.json}; LIBERO-LONG and CALVIN from
  Tables~3--4 of D-LAPA.}
  \label{fig:heatmap}
\end{figure*}

Figure~\ref{fig:interaction} presents a two-panel factorial interaction plot that
disentangles the depth-image effect from the supervision-type effect.
The left panel shows that depth improves performance under both supervision
regimes, but the gain depends strongly on the RGB conditioning strategy:
\mbox{+3.0\,pp} under discrete supervision (\mbox{M3\,$\to$\,M2}) and
\mbox{+6.8\,pp} under continuous distillation (\mbox{M5\,$\to$\,M4}).
The right panel summarises the three pairwise ablation contrasts: the
RGB conditioning effect (\mbox{M1\,$\to$\,M2}, \mbox{+3.4\,pp}) confirms that
switching from discrete index to continuous feature conditioning provides an
additional benefit beyond depth alone, consistent with the paper's conclusion
that continuous RGB features provide richer steering context for the depth encoder.

\begin{figure*}[t]
  \centering
  \includegraphics[width=0.85\textwidth]{fig_interaction.pdf}
  \caption{Left: $2{\times}2$ factorial interaction between depth-image
  presence and supervision type (LIBERO-LONG success rate).
  Right: effect-size bars for each ablation contrast.
  Source: D-LAPA Table~3.}
  \label{fig:interaction}
\end{figure*}

Figure~\ref{fig:calvin} decomposes the CALVIN sequential-task results in two
complementary ways.
The left panel shows per-step $\Delta$ success rate (Model vs.\ LAPA baseline),
revealing a characteristic pattern: both depth models sacrifice approximately
$2{-}3$\,pp on Task~1 but recover strongly at Task~2 ($+4.3$\,pp for Model~2,
$+2.0$\,pp for Model~4) and hold a positive margin through Task~5.
The right panel converts cumulative success rates to conditional transition
probabilities $P(\text{task }k \mid \text{task }k{-}1)$, showing that LAPA's
weakest link is the T4$\to$T5 transition (65.1\%), while both depth models
hold 70--72\% at that step.

\begin{figure*}[t]
  \centering
  \includegraphics[width=0.85\textwidth]{fig_calvin_survival.pdf}
  \caption{CALVIN reanalysis (D-LAPA Table~4).
  Left: per-task $\Delta$ success rate vs.\ LAPA baseline.
  Right: conditional transition probabilities $P(k \mid k{-}1)$.}
  \label{fig:calvin}
\end{figure*}

%%============================================================
\section{Discussion}
%%============================================================

\paragraph{What the reproduction establishes.}
The left-block results (Table~\ref{tab:left}) are numerically faithful within
statistical noise attributable to the split difference.
The right-block results (Table~\ref{tab:right}) reproduce the relative ordering
perfectly under the base LAPA proxy, and the residual analysis adds a
scale-stratified dimension not reported in the paper.
Three independent metrics --- probe $R^2$, Moran's~$I$, and LIBERO/CALVIN
downstream performance --- converge on the same conclusion: depth imagery is
the operative signal.

\paragraph{What could not be reproduced.}
Three items remain out of reach from the provided cached features.
\begin{enumerate}[noitemsep]
  \item The finetuned LAPA right-block baseline ($R^2{=}0.619$), which requires
        the action-supervised LAPA checkpoint. Its absence shifts the right-block
        absolute values by approximately 0.169 $R^2$ points --- not a rounding
        error. This is the single most significant reproducibility gap.
  \item Exact Moran's~$I$ magnitudes, which inherit the same limitation.
  \item The data-size ablation curves from Supplemental~\S.2, which require
        intermediate training checkpoints at 0--65k optimiser steps.
\end{enumerate}

I recommend that the final camera-ready submission archive the finetuned LAPA
features alongside the existing shards.

\paragraph{Limitations.}
The probe target ($|\Delta t|$, a single scalar) does not cover rotational
end-effector components, gripper aperture, or task-semantic structure.
Probe $R^2$ is an in-distribution decodability measure and is not by itself a
predictor of policy success in novel scenes.
The UMAP visualisations are stochastic and non-unique; the chosen
hyperparameters ($n_\text{neighbors}{=}15$, $\text{min\_dist}{=}0.1$) follow
the paper's description.

%%============================================================
\section{Conclusion}
%%============================================================

I have reproduced the feature analysis of D-LAPA within $|\Delta R^2| \le
0.004$ across all reproducible rows and confirmed the paper's central finding
through three independent metrics.
The key result is robust: depth imagery, not model capacity, drives the
geometric grounding improvements in the Stage-2.5 depth encoder family.
The novel scale-stratified residual analysis provides a finer-grained
characterisation: Model~2 achieves an 11.4\% large-displacement MAE reduction,
while identical-architecture no-depth ablations (Models~3 and~5) stay within
$\pm 0.8\%$ of the base LAPA baseline at all displacement scales.
Three novel figures further consolidate the paper's multi-table results into
a form directly suitable for an ICRA submission.

\end{document}